% =====================================================================
\section*{Contributions}\label{sec:contrib}
\addcontentsline{toc}{section}{Contributions}

The work was done by two people in two lanes with one shared pre-registration and one shared decision log. Everything listed under one author was done on that author's machine; the two lanes exchanged only numbered hand-offs. The hand-off codes used in this report are: D2, reference tensors for the merge-parity check; D4, the circuit tensor map; D5b, the steering vector exported as a llama.cpp control vector; D6, the Mac lane's answer files; D7, the central labels returned to the Mac lane; D8, the Mac lane's results and figures.

\paragraph{Sushovit Nanda (first author; project lead; laptop lane).} Led the project from start to finish: conceived and designed the study, set its six questions, its pre-registration and its decision rules, directed both lanes and all of the AI-assisted engineering, and owns every judgement call and every reported result. Wrote Plan v3, the pre-registration and its amendments, and the decision log. Built the data sets (the Aya, XSTest and MedMCQA splits, the ClinicSafe-Hi benchmark and its Hindi and Hinglish variants, the token-matched control corpus) and the laptop harness. Ran the fine-tunes, the merges, the ReSta and circuit-restricted ReSta sweeps, the steering sweep, the NF4 cells and every laptop test cell. Ran the circuit discovery, the necessity and sufficiency validation and the overlap analysis, and produced the tensor map (D4). Wrote the specifications and the style guide from which the ClinicSafe-Hi items and the Hindi and Hinglish variants were drafted, and reviewed the drafts. Designed the hand-label rubric, drew the 300-item sample, and checked and approved every hand label (with the second author); he is the point of contact for every label and every annotation decision in this report. Holds sole authority over, and responsibility for, every judgement call in the decision log: he made the early ones directly (judge choice per language, the stop-token amendment), and the later ones (the trade-off definition, the switch to steering, the handling of degenerate output and of the Hinglish RQ6 contrasts) were proposed by an AI assistant under his written delegation and ratified by him; no decision stands that he has not reviewed. Ran all judging centrally, on one machine, for every cell from both lanes, including the back-translation pipeline, and returned the labels to the Mac lane (D7). Exported the steering vector as a control vector (D5b) and the reference tensors (D2). Ran the statistics, the degeneracy and non-degenerate analyses, the judge validation, the audits and the majority votes, validated and signed off every reported result, and wrote this report.

\paragraph{Manish Shaw (second author; Mac lane).} Built the GGUF conversion and quantization pipeline, the llama.cpp serving and generation harness, the template and EOS audits, and the Mac pre-registrations. Converted and quantized the models to F16, Q4\_K\_M and Q3\_K\_M, ran the fidelity gates and every Mac deployment cell (the unsteered medical cells, the steered cells at five settings and the protected and random Q3\_K\_M cells). Designed and built the equal-size protected and random builds, including the corrected version that excludes the layer-0 tensors, with the size, type and byte-identity checks and the operational checks. Measured the numeric differences between the two lanes' merges (D2). Independently reproduced the deployment and RQ6 statistics from the central labels, computed the conservative union Holm family, ran the throughput and conditional cost analysis, prepared the model card, licence check and redacted submission bundle, and wrote the Mac-lane reports (D8).

\paragraph{Shared.} Both authors reviewed and approved the hand labels, read each other's lane and wrote independent reviews of it, and agreed the amendments. Decision-making and the validation of results sat with the first author throughout. The second and third readings, the adjudication and the blind audits were run by AI agents under his instructions and are reported as AI readings, with the restrictions stated in Section~\ref{sec:judges}; no result in this report rests on a label that was not reviewed by both authors.

\section*{Artifacts}
\addcontentsline{toc}{section}{Artifacts}
The submission consists of the following public artifacts.
\begin{itemize}\setlength{\itemsep}{0pt}
\item \textbf{Hugging Face model collection} (public, un-gated), \url{https://huggingface.co/collections/SushiOrbit/safe-from-fine-tune-to-phone-models}: the LoRA adapters, the merged FP16 weights, the F16 GGUF with the quantization rules and checksums that define the evaluated Q3\_K\_M builds, and the refusal steering vector, each with a model card, under CC-BY-NC-4.0.
\item \textbf{Hugging Face dataset collection} (public), \url{https://huggingface.co/collections/SushiOrbit/safe-from-fine-tune-to-phone-datasets}: ClinicSafe-Hi, the evaluation, calibration and circuit sets, the fine-tuning sets, and the evaluation logs and judge labels, each with a dataset card, under CC-BY-NC-4.0.
\item \textbf{GitHub repository}, \url{https://github.com/SushovitNanda/Safe-from-Fine-Tune-to-Phone}: the code of both lanes, the data package, the results, the pre-registration and decision log, and the reproduction instructions in its README.
\item No demo Space, video or slides accompany this submission.
\end{itemize}

\section*{Repository Layout and Reproduction}
\addcontentsline{toc}{section}{Repository Layout and Reproduction}
The submission repository is organised by stage: a data folder with every evaluation and fine-tuning set as prompt files with a manifest and checksums, and the scripts that build them; a fine-tuning folder with the LoRA training and merge scripts; the laptop harness (cell runner, the two judges, back-translation, statistics and $\kappa$); a repair folder with ReSta, its equivalence check and the steering calibration; a locate folder with the CircuitKIT runs and the circuit validation; the Mac lane's conversion, serving, protection and statistics scripts; a results folder with text-free records, labels and a run record per cell; and the pre-registration, the decision log and the pinned environment at the top level. Reproduction, seed 13 throughout: download the models and sources; build the sets with the data scripts in the documented order, ending with the curation and leak check; train the two LoRAs and merge them; run a cell, judge it (back-translation then WildGuard for Hindi), then run the statistics; run circuit discovery and validation; calibrate the steering vector and export it as a control vector; on the Mac, serve a GGUF with the pinned llama.cpp build using the settings in Section~\ref{sec:compute} and drive it with the Mac cell runner. Environment: Python 3.12.10, torch 2.11.0 (CUDA 12.8), transformers 5.22.1, peft 0.21.2, bitsandbytes 0.50.2, curatorkit 0.1.21, safetune 0.1.8, circuitkit 0.2.0, transformer\_lens 3.8.0, auditkit 1.1.3; Mac dependencies pinned in the repository.

\section*{Data and Code Availability}
\addcontentsline{toc}{section}{Data and Code Availability}
The submission repository contains the harness, scripts, configurations, pre-registration, decision log, text-free per-answer records, judge labels, derived statistics and this report. It omits raw model answers to harmful prompts, back-translated text, reference tensors and access tokens, which stay in the authors' private working repository; model weights are on Hugging Face. The ClinicSafe-Hi prompts are released in clear text with their labels; only model completions are withheld. Tiny Aya Global (one pinned snapshot) is used under its CC-BY-NC-4.0 licence, with Tiny Aya Base as the ReSta base; Aya and MedMCQA are Apache-2.0, XSTest CC-BY-4.0, Alpaca-cleaned CC-BY-NC-4.0, and ClinicSafe-Hi is released under CC-BY-NC-4.0. Weights, the control vector and the adapters are redistributed only through the model collection above.

\paragraph{Work done outside the hackathon.} None. The project's working repository (private; its whitelist export is the public submission repository) begins on 3 November 2026 with environment setup, as does the decision log; every dataset, script, fine-tune, cell, label and figure, and this report, were produced between 3 and 15 November 2026 by the two authors. The only pre-existing code is third-party: the track toolkits at pinned versions, llama.cpp at a pinned build, TransformerLens 3.8.0, and the public models and benchmarks cited above.

\paragraph{Use of AI assistants.} AI assistants (Claude Code) did the engineering: the complete code base and the environment set-up on both machines, the execution of the fine-tunes, sweeps and evaluation cells, debugging, and the monitoring of logs, all under the authors' written specifications and recorded in the decision log. Decision-making and the validation of results stayed with the first author: he wrote the specifications and the style guide, designed the labelling rubric and drew the sample, checked and approved every hand label with the second author, took or ratified every judgement call, and signed off every reported number. Within that frame, AI agents produced first drafts to his specifications, each reviewed and approved before use (the ClinicSafe-Hi items, the Hindi and Hinglish variants and post-edits, and the 300 hand labels); the blind second and third readings, the adjudication, the audits and the majority votes, which are reported as AI readings; and the drafts of Amendments 3 to 5, which he ratified. A large share of the code and of this text was drafted with AI assistance and edited by the authors. No number in this report was written by an assistant without a script reading it from the result files, and the independent re-derivation in Section~\ref{sec:holm} confirms them.

\paragraph{Changelog.} 3 Nov: environment, data package v1, Plan v2. 5 Nov: Plan v3 (six questions) and the pre-registration; Amendment 1. 8 Nov: Amendment 2. 9 Nov: Amendments 3 and 4; Mac protocol. 11 Nov: Amendment 5 written and withdrawn the same day; corrected protected and random builds registered and run; RQ6 judged; first full build of this report. 13 Nov: Mac D8 integrated; all results frozen. 14 Nov: executive summary, hypotheses, artifacts, reproduction and this changelog added after three rounds of independent review. 15 Nov: artifact links, wording and consistency corrections after a final five-way review; public release of the artifacts. Later on 15 Nov: a two-series headline figure (all languages and English plus Hindi), the English-plus-Hindi result with its intervals in the executive summary, bootstrap intervals and item counts for every $\kappa$, and larger table type.

\section*{Ethics Statement}
\addcontentsline{toc}{section}{Ethics Statement}
The harmful prompts come from existing safety benchmarks or were drafted by AI agents to the authors' specification and reviewed by them; no harmful model output is quoted in this report or in the submission repository. The fine-tuned weights are released through the collection described under Artifacts (merged FP16 weights and the F16 GGUF, with the rules and checksums that define the quantized builds) together with the control vector, so that every result can be reproduced; the full-precision weights add no risk beyond the quantized medical model, and the drift they carry is the finding itself. \textbf{Responsible disclosure.} The findings concern a public model under standard fine-tuning and serving settings, not a vulnerability in any third-party system; no raw harmful completions are released; the organisers receive the full results through the submission, and Cohere Labs can be given the per-language rates on request. The clinical prompts are synthetic and contain no patient data. The authors see a dual-use risk in publishing which deployment settings weaken refusals; we judge that risk small, because the settings are standard and the effects are modest, and outweighed by the value to deployers of knowing them.

\section*{Acknowledgements}
\addcontentsline{toc}{section}{Acknowledgements}
We thank the organisers of the Lexsi AI Safety and Alignment Hackathon for the Track 2 brief, and Cohere Labs for releasing Tiny Aya. All compute was the authors' own hardware, apart from one parity check on a free Colab T4; the cash cost of compute was zero.
